% !TeX program = lualatex
% !TeX root = pred2.tex
% ============================================================
% Zajednički preamble za sva predavanja RWA kolegija
% Uključiti s: % ============================================================
% Zajednički preamble za sva predavanja RWA kolegija
% Uključiti s: % ============================================================
% Zajednički preamble za sva predavanja RWA kolegija
% Uključiti s: \input{../common/preamble.tex}
% ============================================================

\documentclass[aspectratio=149,11pt]{beamer}
\usetheme{CambridgeUS}

\usepackage{fontspec}              % radi samo sa XeLaTeX/LuaLaTeX
\setsansfont{Latin Modern Sans}    % sans-serif font
\setmonofont{Latin Modern Mono}    % monospace font za kod
\usepackage[croatian]{babel}
\usepackage{csquotes}              % \enquote{} → hrvatski navodnici „..."
\usepackage{lmodern}
\usepackage{booktabs}
\usepackage{listings}
\usepackage{xcolor}
\usepackage{hyperref}

% ---------- Autor i logo ----------
\author{mag.ing. Josip Torić (SIT UNIZD)}
\titlegraphic{\includegraphics[height=2.5cm]{../common/SIT_logo_tekst1.pdf}}

% ---------- Listings ----------
\definecolor{codebg}{RGB}{248,248,248}
\definecolor{codeframeclr}{RGB}{220,220,220}
\lstdefinelanguage{bash}{
  morekeywords={sudo,apt,brew,pip,pipx,uv,winget,python,git,flask,source,deactivate,code,gunicorn,docker,compose,heroku,render,railway,node,npm,pnpm,yarn,uvicorn,alembic,pytest,ruff,black,precommit},
  sensitive=false,
  morecomment=[l]{\#}
}
\lstdefinelanguage{python}{
  morekeywords={from,import,def,class,return,if,elif,else,for,in,while,try,except,raise,with,as,not,and,or,is,lambda,yield,True,False,None,await,async},
  sensitive=true,
  morecomment=[l]{\#},
  morestring=[b]',
  morestring=[b]"
}
\lstdefinelanguage{sql}{
  morekeywords={SELECT,FROM,WHERE,INSERT,INTO,VALUES,UPDATE,SET,DELETE,CREATE,TABLE,PRIMARY,KEY,FOREIGN,REFERENCES,UNIQUE,NOT,NULL,INDEX,ON,ALTER,DROP,CONSTRAINT},
  sensitive=false,
  morecomment=[l]{--}
}
\lstdefinelanguage{json}{
  morecomment=[l]{//},
  morestring=[b]",
  showstringspaces=false
}
\lstdefinelanguage{typescript}{
  morekeywords={import,export,from,const,let,var,function,return,if,else,for,of,in,while,switch,case,break,async,await,try,catch,finally,throw,new,class,this,extends,implements,interface,type,enum,readonly,public,private,protected,static,get,set,true,false,null,undefined,void,string,number,boolean,any,unknown,never,typeof,keyof,as},
  sensitive=true,
  morecomment=[l]{//},
  morecomment=[s]{/*}{*/},
  morestring=[b]',
  morestring=[b]",
  morestring=[b]`
}
\lstalias{javascript}{typescript}
\lstset{
  backgroundcolor=\color{codebg},
  basicstyle=\ttfamily\footnotesize,
  frame=single,
  rulecolor=\color{codeframeclr},
  frameround=tttt,
  breaklines=true,
  tabsize=2,
  showstringspaces=false,
  upquote=true,
  columns=fullflexible
}

% ---------- Helper macros ----------
\newcommand{\codefile}[1]{\texttt{#1}}
\newcommand{\kbd}[1]{\texttt{#1}}

% ---------- Code frame environment ----------
% Automatski dodaje fragile i shrink=5, koristi se umjesto \begin{frame}
% Primjer: \begin{codeframe}{Naslov} ... \end{codeframe}
% Primjer s opcijama: \begin{codeframe}[shrink=15]{Naslov} ... \end{codeframe}
\newenvironment{codeframe}[2][]
  {\begin{frame}[fragile,shrink=5,#1]{#2}}
  {\end{frame}}

% ============================================================

\documentclass[aspectratio=149,11pt]{beamer}
\usetheme{CambridgeUS}

\usepackage{fontspec}              % radi samo sa XeLaTeX/LuaLaTeX
\setsansfont{Latin Modern Sans}    % sans-serif font
\setmonofont{Latin Modern Mono}    % monospace font za kod
\usepackage[croatian]{babel}
\usepackage{csquotes}              % \enquote{} → hrvatski navodnici „..."
\usepackage{lmodern}
\usepackage{booktabs}
\usepackage{listings}
\usepackage{xcolor}
\usepackage{hyperref}

% ---------- Autor i logo ----------
\author{mag.ing. Josip Torić (SIT UNIZD)}
\titlegraphic{\includegraphics[height=2.5cm]{../common/SIT_logo_tekst1.pdf}}

% ---------- Listings ----------
\definecolor{codebg}{RGB}{248,248,248}
\definecolor{codeframeclr}{RGB}{220,220,220}
\lstdefinelanguage{bash}{
  morekeywords={sudo,apt,brew,pip,pipx,uv,winget,python,git,flask,source,deactivate,code,gunicorn,docker,compose,heroku,render,railway,node,npm,pnpm,yarn,uvicorn,alembic,pytest,ruff,black,precommit},
  sensitive=false,
  morecomment=[l]{\#}
}
\lstdefinelanguage{python}{
  morekeywords={from,import,def,class,return,if,elif,else,for,in,while,try,except,raise,with,as,not,and,or,is,lambda,yield,True,False,None,await,async},
  sensitive=true,
  morecomment=[l]{\#},
  morestring=[b]',
  morestring=[b]"
}
\lstdefinelanguage{sql}{
  morekeywords={SELECT,FROM,WHERE,INSERT,INTO,VALUES,UPDATE,SET,DELETE,CREATE,TABLE,PRIMARY,KEY,FOREIGN,REFERENCES,UNIQUE,NOT,NULL,INDEX,ON,ALTER,DROP,CONSTRAINT},
  sensitive=false,
  morecomment=[l]{--}
}
\lstdefinelanguage{json}{
  morecomment=[l]{//},
  morestring=[b]",
  showstringspaces=false
}
\lstdefinelanguage{typescript}{
  morekeywords={import,export,from,const,let,var,function,return,if,else,for,of,in,while,switch,case,break,async,await,try,catch,finally,throw,new,class,this,extends,implements,interface,type,enum,readonly,public,private,protected,static,get,set,true,false,null,undefined,void,string,number,boolean,any,unknown,never,typeof,keyof,as},
  sensitive=true,
  morecomment=[l]{//},
  morecomment=[s]{/*}{*/},
  morestring=[b]',
  morestring=[b]",
  morestring=[b]`
}
\lstalias{javascript}{typescript}
\lstset{
  backgroundcolor=\color{codebg},
  basicstyle=\ttfamily\footnotesize,
  frame=single,
  rulecolor=\color{codeframeclr},
  frameround=tttt,
  breaklines=true,
  tabsize=2,
  showstringspaces=false,
  upquote=true,
  columns=fullflexible
}

% ---------- Helper macros ----------
\newcommand{\codefile}[1]{\texttt{#1}}
\newcommand{\kbd}[1]{\texttt{#1}}

% ---------- Code frame environment ----------
% Automatski dodaje fragile i shrink=5, koristi se umjesto \begin{frame}
% Primjer: \begin{codeframe}{Naslov} ... \end{codeframe}
% Primjer s opcijama: \begin{codeframe}[shrink=15]{Naslov} ... \end{codeframe}
\newenvironment{codeframe}[2][]
  {\begin{frame}[fragile,shrink=5,#1]{#2}}
  {\end{frame}}

% ============================================================

\documentclass[aspectratio=149,11pt]{beamer}
\usetheme{CambridgeUS}

\usepackage{fontspec}              % radi samo sa XeLaTeX/LuaLaTeX
\setsansfont{Latin Modern Sans}    % sans-serif font
\setmonofont{Latin Modern Mono}    % monospace font za kod
\usepackage[croatian]{babel}
\usepackage{csquotes}              % \enquote{} → hrvatski navodnici „..."
\usepackage{lmodern}
\usepackage{booktabs}
\usepackage{listings}
\usepackage{xcolor}
\usepackage{hyperref}

% ---------- Autor i logo ----------
\author{mag.ing. Josip Torić (SIT UNIZD)}
\titlegraphic{\includegraphics[height=2.5cm]{../common/SIT_logo_tekst1.pdf}}

% ---------- Listings ----------
\definecolor{codebg}{RGB}{248,248,248}
\definecolor{codeframeclr}{RGB}{220,220,220}
\lstdefinelanguage{bash}{
  morekeywords={sudo,apt,brew,pip,pipx,uv,winget,python,git,flask,source,deactivate,code,gunicorn,docker,compose,heroku,render,railway,node,npm,pnpm,yarn,uvicorn,alembic,pytest,ruff,black,precommit},
  sensitive=false,
  morecomment=[l]{\#}
}
\lstdefinelanguage{python}{
  morekeywords={from,import,def,class,return,if,elif,else,for,in,while,try,except,raise,with,as,not,and,or,is,lambda,yield,True,False,None,await,async},
  sensitive=true,
  morecomment=[l]{\#},
  morestring=[b]',
  morestring=[b]"
}
\lstdefinelanguage{sql}{
  morekeywords={SELECT,FROM,WHERE,INSERT,INTO,VALUES,UPDATE,SET,DELETE,CREATE,TABLE,PRIMARY,KEY,FOREIGN,REFERENCES,UNIQUE,NOT,NULL,INDEX,ON,ALTER,DROP,CONSTRAINT},
  sensitive=false,
  morecomment=[l]{--}
}
\lstdefinelanguage{json}{
  morecomment=[l]{//},
  morestring=[b]",
  showstringspaces=false
}
\lstdefinelanguage{typescript}{
  morekeywords={import,export,from,const,let,var,function,return,if,else,for,of,in,while,switch,case,break,async,await,try,catch,finally,throw,new,class,this,extends,implements,interface,type,enum,readonly,public,private,protected,static,get,set,true,false,null,undefined,void,string,number,boolean,any,unknown,never,typeof,keyof,as},
  sensitive=true,
  morecomment=[l]{//},
  morecomment=[s]{/*}{*/},
  morestring=[b]',
  morestring=[b]",
  morestring=[b]`
}
\lstalias{javascript}{typescript}
\lstset{
  backgroundcolor=\color{codebg},
  basicstyle=\ttfamily\footnotesize,
  frame=single,
  rulecolor=\color{codeframeclr},
  frameround=tttt,
  breaklines=true,
  tabsize=2,
  showstringspaces=false,
  upquote=true,
  columns=fullflexible
}

% ---------- Helper macros ----------
\newcommand{\codefile}[1]{\texttt{#1}}
\newcommand{\kbd}[1]{\texttt{#1}}

% ---------- Code frame environment ----------
% Automatski dodaje fragile i shrink=5, koristi se umjesto \begin{frame}
% Primjer: \begin{codeframe}{Naslov} ... \end{codeframe}
% Primjer s opcijama: \begin{codeframe}[shrink=15]{Naslov} ... \end{codeframe}
\newenvironment{codeframe}[2][]
  {\begin{frame}[fragile,shrink=5,#1]{#2}}
  {\end{frame}}


\title[Razvoj web aplikacija]{Razvoj web aplikacija - Predavanje 2 \\ Modeliranje podataka, ORM i migracije}
\date{12.10.2026.}

\begin{document}
	
	\frame{\titlepage}
	
	% ------------------------------------------------------------
	\begin{frame}{Plan predavanja}
		\begin{enumerate}
			\item Modeliranje podataka: entiteti, relacije, constrainti
			\item Normalizacija (minimalno) i tipične greške u modelu
			\item SQL kao temelj: što ORM zapravo radi
			\item SQLAlchemy 2.0: modeli, engine, session
			\item Migracije: Alembic, autogenerate, rollback
			\item Seed podaci i razvojni workflow
		\end{enumerate}
		
		\vspace{0.4em}
		\begin{block}{Cilj}
			Na kraju predavanja imamo bazu s tablicama \textbf{Club} i \textbf{User}, migracije i seed.
		\end{block}
	\end{frame}
	
	% ------------------------------------------------------------
	\begin{frame}{Ishodi učenja}
		Nakon ovog predavanja student može:
		\begin{itemize}
			\item modelirati tablice i relacije za realnu web aplikaciju
			\item objasniti razliku: model (DB) vs schema (API DTO)
			\item postaviti SQLAlchemy engine i session per request
			\item koristiti Alembic za izradu i primjenu migracija
			\item objasniti unique/foreign key constraint i zašto su bitni
		\end{itemize}
	\end{frame}
	
	% ------------------------------------------------------------
	\begin{frame}{Prvo: domena našeg projekta (što modeliramo?)}
		Minimalni entiteti za početak:
		\begin{itemize}
			\item \textbf{Club} - klub (naziv, grad, kontakt email, kontakt telefon)
			\item \textbf{User} - korisnik (admin ili club user)
		\end{itemize}
		
		\vspace{0.6em}
		Kasnije dodajemo:
		\begin{itemize}
			\item \textbf{Lifter} (natjecatelj), \textbf{Competition}, \textbf{Registration}
		\end{itemize}
		
		\vspace{0.6em}
		Iterativno: prvo temelj (auth + ownership), tek onda složenije entitete.
	\end{frame}
	
	% ------------------------------------------------------------
	\begin{frame}{Entitet, tablica, redak (terminologija)}
		\begin{itemize}
			\item \textbf{Entitet} (concept) $\rightarrow$ \textbf{tablica} (implementation)
			\item \textbf{Instanca entiteta} $\rightarrow$ \textbf{redak} (row)
			\item \textbf{Atribut entiteta} $\rightarrow$ \textbf{stupac} (column)
		\end{itemize}
		
		\vspace{0.6em}
		Modeliramo \textbf{stabilne} pojmove domene, ne UI forme.
	\end{frame}
	
	% ------------------------------------------------------------
	\begin{frame}{Primarni i strani ključ}
		\begin{columns}[T]
			\begin{column}{0.48\textwidth}
				\textbf{Primarni ključ (PK)}
				\begin{itemize}
					\item jednoznačno identificira redak
					\item najčešće \textbf{surrogate key}: \texttt{id} integer/uuid
					\item prirodni ključ (OIB, registracijski broj) - oprez
				\end{itemize}
			\end{column}
			\begin{column}{0.48\textwidth}
				\textbf{Strani ključ (FK)}
				\begin{itemize}
					\item povezuje dvije tablice
					\item referencijalni integritet: nema usera s \texttt{club\_id} koji ne postoji
					\item određuje relaciju (1:N, M:N)
				\end{itemize}
			\end{column}
		\end{columns}

		\vspace{0.8em}
		Prirodni podaci se mijenjaju ili nisu uvijek dostupni, \texttt{id} je stabilan.
		U našem projektu: Club 1\,:\,1 User, Club 1\,:\,N Lifter.
	\end{frame}
	
	% ------------------------------------------------------------
	\begin{frame}{Relacije: 1:1, 1:N, M:N (brzi pregled)}
		\begin{center}
			\small
			\begin{tabular}{lll}
				\toprule
				\textbf{Relacija} & \textbf{Primjer iz projekta} & \textbf{Implementacija} \\
				\midrule
				1\,:\,1 & Club $\leftrightarrow$ User & FK + UNIQUE; rijetko, često ista tablica \\
				1\,:\,N & Club $\rightarrow$ Lifter & FK na strani \enquote{N}; najčešća relacija \\
				M\,:\,N & Lifter $\leftrightarrow$ Competition & pomoćna tablica (Registration) \\
				\bottomrule
			\end{tabular}
		\end{center}

		\vspace{0.6em}
		Kad dođemo do prijava: Competition 1\,:\,N Registration i Lifter 1\,:\,N Registration.
	\end{frame}
	
	% ------------------------------------------------------------
	\begin{frame}{Constrainti: što baza smije/ne smije dopustiti}
		\begin{center}
			\small
			\begin{tabular}{lll}
				\toprule
				\textbf{Constraint} & \textbf{Što jamči} & \textbf{Primjer} \\
				\midrule
				NOT NULL & vrijednost mora postojati & \texttt{name} kluba \\
				UNIQUE & nema duplikata & \texttt{email}, \texttt{username} \\
				CHECK & vrijednost poštuje pravilo & težina $>$ 0 \\
				FOREIGN KEY & relacija mora biti valjana & \texttt{club\_id} postoji u \texttt{clubs} \\
				\bottomrule
			\end{tabular}
		\end{center}
	\end{frame}

	% ------------------------------------------------------------
	\begin{frame}{}
		\begin{center}
			{\Large Pravila u kodu nisu dovoljna.}

			\vspace{1em}
			{\huge \textbf{Baza je autoritet.}}
		\end{center}
	\end{frame}
	
	% ------------------------------------------------------------
	\begin{frame}{Unique constraint: primjer s usernameom}
		\begin{itemize}
			\item Username korisnika mora biti jedinstven
			\item Bez UNIQUE: login postaje dvosmislen ili nesiguran
			\item Unique constraint je zaštita od race conditiona:
			\begin{itemize}
				\item dva requesta u isto vrijeme pokušaju kreirati isti username
			\end{itemize}
		\end{itemize}
		
		\vspace{0.6em}
		Na obrani pitam: \emph{gdje} se jedinstvenost provjerava - u kodu, u bazi ili oboje?
	\end{frame}
	
	% ------------------------------------------------------------
	\begin{frame}{Normalizacija (minimalno što trebate znati)}
		\begin{itemize}
			\item Cilj: smanjiti dupliciranje i anomalije (update/insert/delete)
			\item 1NF: atomarne vrijednosti (nema lista u stupcu)
			\item 2NF/3NF: atributi ovise o ključu, a ne o drugim atributima
		\end{itemize}
		
		\vspace{0.6em}
		Ako imate isti podatak na 3 mjesta, negdje će kad-tad postati netočan.
	\end{frame}
	
	% ------------------------------------------------------------
	\begin{frame}{Tipične greške u modeliranju (iz prakse)}
		\begin{center}
			\small
			\begin{tabular}{ll}
				\toprule
				\textbf{Greška} & \textbf{Posljedica} \\
				\midrule
				lista u jednom stupcu (\texttt{"a,b,c"}) & nema query-ja po elementima \\
				email kao primarni ključ & email se može promijeniti \\
				bez constrainta \enquote{jer ćemo provjeriti u kodu} & race condition \\
				denormalizacija bez razloga & neusklađenost podataka \\
				\bottomrule
			\end{tabular}
		\end{center}

		\vspace{0.6em}
		Ispravan model = manje bugova i manje if-ova u kodu.
	\end{frame}
	
	% ------------------------------------------------------------
	\begin{frame}{SQL je temelj (ORM je sloj iznad)}
		\begin{itemize}
			\item ORM generira SQL (ili ga mapira)
			\item Ako ne razumijete SQL:
			\begin{itemize}
				\item ne znate zašto je query spor
				\item ne znate zašto relationship radi čudno
			\end{itemize}
			\item Danas radimo ORM, ali s mentalnim modelom SQL-a
		\end{itemize}
	\end{frame}
	
	% ------------------------------------------------------------
	\begin{frame}[fragile,shrink=5]{SQL: minimalni CRUD primjeri}
\begin{lstlisting}[language=sql]
-- CREATE
INSERT INTO clubs (name, city) VALUES ('Behemot', 'Zadar');

-- READ
SELECT id, name FROM clubs WHERE city = 'Zadar';

-- UPDATE
UPDATE clubs SET city = 'Biograd' WHERE id = 1;

-- DELETE
DELETE FROM clubs WHERE id = 1;
\end{lstlisting}
	\end{frame}
	
	% ------------------------------------------------------------
	\begin{frame}[fragile,shrink=5]{JOIN: zašto relacije postoje}
\begin{lstlisting}[language=sql]
SELECT u.id, u.username, c.name AS club_name
FROM users u
JOIN clubs c ON c.id = u.club_id
WHERE u.role = 'club';
\end{lstlisting}
		
		\vspace{0.6em}
		Relacije omogućuju query-je koji su prirodni i učinkoviti.
	\end{frame}
	
	% ------------------------------------------------------------
	\begin{frame}{Index: kada i zašto?}
		\begin{itemize}
			\item Index ubrzava \texttt{WHERE}, \texttt{JOIN}, \texttt{ORDER BY}
			\item Index košta:
			\begin{itemize}
				\item zauzima prostor
				\item usporava INSERT/UPDATE (mora se ažurirati)
			\end{itemize}
			\item Unique constraint obično stvara index (npr. email)
		\end{itemize}
		
		\vspace{0.6em}
		Indeksiraj ono po čemu često filtriraš: \texttt{email}, \texttt{club\_id}, \texttt{competition\_id}.
	\end{frame}
	
	% ------------------------------------------------------------
	\begin{frame}{Transakcije (što su i zašto?)}
		\begin{itemize}
			\item Transakcija = skup DB operacija koje su \textbf{atomarne}
			\item ACID (kratko):
			\begin{itemize}
				\item Atomicity: sve ili ništa
				\item Consistency: constrainti vrijede
				\item Isolation: konkurentni requestovi
				\item Durability: zapis ostaje
			\end{itemize}
		\end{itemize}
		
		\vspace{0.6em}
		Svaki request tipično ima svoju transakciju (commit/rollback).
	\end{frame}
	
	% ------------------------------------------------------------
	\begin{frame}{Model vs Schema (DB vs API)}
		\begin{itemize}
			\item \textbf{Model} (SQLAlchemy): kako se podatak sprema u bazu
			\item \textbf{Schema} (Pydantic): kako podatak putuje preko API-ja
		\end{itemize}
		
		\vspace{0.6em}
		Zašto razdvajamo?
		\begin{itemize}
			\item ne želimo izlagati sve DB stupce (npr. \texttt{password\_hash})
			\item različiti input/output oblici (create vs read)
			\item sigurnost i kontrola ugovora
		\end{itemize}
	\end{frame}
	
	% ------------------------------------------------------------
	\begin{frame}{SQLAlchemy 2.0: osnovni komadi}
		\begin{itemize}
			\item \textbf{Engine}: konekcija prema DB-u
			\item \textbf{Session}: jedinica rada (unit of work)
			\item \textbf{Declarative model}: Python klase kao tablice
			\item \textbf{Query}: SELECT/INSERT/UPDATE/DELETE kroz session
		\end{itemize}
		
		\vspace{0.6em}
		Danas postavljamo engine + session i kreiramo prva dva modela.
	\end{frame}
	
	% ------------------------------------------------------------
	\begin{frame}[fragile, shrink=15]{Engine i session (async)}
\begin{lstlisting}[language=python]
from sqlalchemy.ext.asyncio import (
    create_async_engine, async_sessionmaker, AsyncSession
)

engine = create_async_engine(DATABASE_URL, pool_pre_ping=True)
AsyncSessionLocal = async_sessionmaker(
    bind=engine, class_=AsyncSession,
    expire_on_commit=False
)

async def get_db():
    async with AsyncSessionLocal() as session:
        try:
            yield session
            await session.commit()
        except:
            await session.rollback()
            raise
\end{lstlisting}

		\vspace{0.4em}
		Async session per request, commit/rollback kontrolirano. Koristimo \texttt{asyncpg} driver.
		Connection string za async mora biti \texttt{postgresql+asyncpg://...}, a ne obični \texttt{postgresql://}.
		To je čest uzrok da projekt \enquote{radi na predavanju, a ne radi doma}.
	\end{frame}
	
	% ------------------------------------------------------------
	\begin{frame}{Zašto session per request?}
		\begin{itemize}
			\item Web aplikacija ima mnogo paralelnih requestova
			\item Jedna globalna session:
			\begin{itemize}
				\item race condition
				\item curenje stanja između requestova
				\item čudni bugovi koje ne možeš reproducirati
			\end{itemize}
		\end{itemize}
		
		\vspace{0.6em}
		U FastAPI dependencyju otvori session, na kraju je zatvori.
	\end{frame}
	
	% ------------------------------------------------------------
	\begin{frame}[fragile,shrink=5]{Model: \texttt{Club} (SQLAlchemy 2.0 style)}
\begin{lstlisting}[language=python]
from sqlalchemy.orm import Mapped, mapped_column, relationship
from sqlalchemy import String

class Club(Base):
    __tablename__ = "clubs"

    id: Mapped[int] = mapped_column(primary_key=True)
    name: Mapped[str] = mapped_column(String(120), unique=True, nullable=False)
    city: Mapped[str] = mapped_column(String(80), nullable=False)

    users: Mapped[list["User"]] = relationship(back_populates="club")
\end{lstlisting}
	\end{frame}
	
	% ------------------------------------------------------------
	\begin{frame}[fragile,shrink=5]{Model: \texttt{User} (role + club\_id)}
\begin{lstlisting}[language=python]
from sqlalchemy import String, ForeignKey
from sqlalchemy.orm import Mapped, mapped_column, relationship

class User(Base):
    __tablename__ = "users"

    id: Mapped[int] = mapped_column(primary_key=True)
    username: Mapped[str] = mapped_column(String(100), unique=True, nullable=False)
    password_hash: Mapped[str] = mapped_column(String(255), nullable=False)
    role: Mapped[str] = mapped_column(String(20), nullable=False)

    club_id: Mapped[int | None] = mapped_column(ForeignKey("clubs.id"), nullable=True)
    club: Mapped["Club | None"] = relationship(back_populates="users")
\end{lstlisting}
		
		\vspace{0.4em}
		Admin može imati \texttt{club\_id = NULL}.
	\end{frame}
	
	% ------------------------------------------------------------
	\begin{frame}{Enum ili string za role?}
		\begin{itemize}
			\item String je jednostavan, ali riskira typo (\enquote{admn})
			\item Enum je sigurniji, ali traži malo više konfiguracije
		\end{itemize}
		
		\vspace{0.6em}
		Možemo krenuti sa stringom + validation u schemi, a kasnije prebaciti u Enum.
	\end{frame}
	
	% ------------------------------------------------------------
	\begin{frame}{Naming conventions (zašto je bitno?)}
		\begin{itemize}
			\item Tablice u množini (\texttt{users}, \texttt{clubs})
			\item Foreign key: \texttt{club\_id}
			\item Index/constraint imena: predvidljiva (pomaže migracijama)
		\end{itemize}
		
		\vspace{0.6em}
		Ako naming nije dosljedan, migracije i debug postaju teži.
	\end{frame}
	
	% ------------------------------------------------------------
	\begin{frame}{Alembic: što je i zašto ga koristimo?}
		\begin{itemize}
			\item Version control za DB shemu
			\item Svaka promjena modela $\rightarrow$ nova migracija
			\item Omogućuje:
			\begin{itemize}
				\item timski rad bez \enquote{radi mi na mojoj bazi}
				\item kontrolirane promjene u produkciji
				\item rollback kad nešto pođe krivo
			\end{itemize}
		\end{itemize}
	\end{frame}
	
	% ------------------------------------------------------------
	\begin{frame}[fragile,shrink=5]{Alembic: osnovne naredbe}
\begin{lstlisting}[language=bash]
# inicijalizacija (jednom)
alembic init alembic

# generiranje migracije iz modela
alembic revision --autogenerate -m "init clubs and users"

# primjena migracija
alembic upgrade head

# rollback
alembic downgrade -1
\end{lstlisting}
	\end{frame}
	
	% ------------------------------------------------------------
	\begin{frame}{Autogenerate nije magija (što provjeriti)}
		\begin{itemize}
			\item Autogenerate uspoređuje \textbf{metadata} s bazom
			\item Ali može pogriješiti ili propustiti:
			\begin{itemize}
				\item složene constraint-e
				\item rename stupaca (može ispasti drop+create)
			\end{itemize}
		\end{itemize}
		
		\vspace{0.6em}
		Uvijek pročitaj migraciju prije \texttt{upgrade}.
	\end{frame}
	
	% ------------------------------------------------------------
	\begin{frame}[fragile,shrink=5]{Kako izgleda migracija (primjer)}
\begin{lstlisting}[language=python]
def upgrade() -> None:
    op.create_table(
        "clubs",
        sa.Column("id", sa.Integer(), primary_key=True),
        sa.Column("name", sa.String(length=120), nullable=False),
        sa.Column("city", sa.String(length=80), nullable=False),
        sa.UniqueConstraint("name"),
    )
\end{lstlisting}
		
		\vspace{0.6em}
		Gdje je UNIQUE i što sprječava?
	\end{frame}
	
	% ------------------------------------------------------------
	\begin{frame}{Rollback (downgrade) i zašto ga trebate znati}
		\begin{itemize}
			\item Kad migracija u produkciji pukne, rollback može spasiti situaciju
			\item U razvoju rollback pomaže za provjeru da migracije nisu \enquote{one-way}
			\item Ali rollback nije uvijek trivijalan - drop stupca znači nepovratan gubitak podataka
		\end{itemize}
		
		\vspace{0.6em}
		Ne oslanjaj se na rollback kao plan. Planiraj migracije.
	\end{frame}
	
	% ------------------------------------------------------------
	\begin{frame}{Seed podaci: što, zašto, kako}
		\begin{itemize}
			\item Seed = inicijalni podaci za razvoj/test (admin, demo klub)
			\item Nije isto što i migracija:
			\begin{itemize}
				\item migracija mijenja shemu
				\item seed puni tablice
			\end{itemize}
			\item Seed mora biti \textbf{idempotentan} ili barem siguran za ponavljanje
		\end{itemize}
	\end{frame}
	
	% ------------------------------------------------------------
	\begin{frame}{Seed: minimalni cilj}
		U seed-u danas želimo:
		\begin{itemize}
			\item kreirati \textbf{admin} korisnika
			\item kreirati jedan \textbf{demo klub}
			\item kreirati jednog \textbf{club usera} za demo klub
		\end{itemize}
		
		\vspace{0.6em}
		Da možemo odmah testirati auth i ownership kada dođemo do JWT predavanja.
	\end{frame}
	
	% ------------------------------------------------------------
	\begin{frame}{Timestampi i timezone (kratko, ali bitno)}
		\begin{itemize}
			\item U produkciji:
			\begin{itemize}
				\item spremamo vrijeme u UTC
				\item prikaz radimo u lokalnoj zoni (UI)
			\end{itemize}
			\item U našem projektu rokovi (prelim/final) će biti datetime:
			\begin{itemize}
				\item backend uspoređuje s server vremenom (UTC)
			\end{itemize}
		\end{itemize}
		
		\vspace{0.6em}
		Na obrani, ako datumi \enquote{ne rade}, pitam: timezone? format? gdje uspoređuješ?
	\end{frame}
	
	% ------------------------------------------------------------
	\begin{frame}{Demo kroz commitove (Predavanje 2)}
		\begin{itemize}
			\item \texttt{feat(api): add SQLAlchemy engine and session factory}
			\item \texttt{feat(api): add Club and User models with relationship}
			\item \texttt{chore(api): add Alembic and initial migration}
			\item \texttt{feat(api): create initial migration for Club and User}
			\item \texttt{feat(api): add seed script for admin and demo club}
		\end{itemize}
		
		\vspace{0.6em}
		\textbf{Na satu:} prolazimo diffeve i provjeravamo DB stanje.
	\end{frame}
	
	% ------------------------------------------------------------
	\begin{frame}{Mini checklista: što mora raditi do kraja sata}
		\begin{enumerate}
			\item \texttt{docker compose up -d db} (Postgres radi)
			\item \texttt{alembic upgrade head} (tablice su kreirane)
			\item seed skripta se izvrši bez greške
			\item u bazi postoji admin user i demo klub
		\end{enumerate}
		
		\vspace{0.6em}
		Ako ovo nemate, sve iduće (JWT, ownership, registracije) će vam biti teže.
	\end{frame}
	
	% ------------------------------------------------------------
	\begin{frame}{Pitanja za diskusiju (za provjeru razumijevanja)}
		\begin{itemize}
			\item Zašto je email UNIQUE? Što se dogodi bez toga?
			\item Zašto admin ima \texttt{club\_id = NULL}?
			\item Što radi foreign key constraint?
			\item Što je razlika model vs schema?
			\item Zašto session per request, a ne globalna session?
		\end{itemize}
	\end{frame}
	
	% ------------------------------------------------------------
	\begin{frame}{Što ide sljedeće (Predavanje 3)}
		\begin{itemize}
			\item Hashiranje lozinki
			\item JWT access + refresh token
			\item Login i refresh endpointi
			\item Dependency \texttt{get\_current\_user}
			\item Sigurnost u SPA (XSS i token storage)
		\end{itemize}
	\end{frame}
	
\end{document}